\documentclass[11pt]{article}
\usepackage{amsmath,amssymb,array,booktabs,geometry,siunitx}
\usepackage[hidelinks]{hyperref}
\geometry{margin=1in}
\begin{document}

\title{Verification, validation, and numerical uncertainty}
\author{cse-heat-equation-example}
\date{}
\maketitle

\section{Evidence layers}

\begin{center}
\begin{tabular}{@{}>{\raggedright\arraybackslash}p{0.22\textwidth}p{0.68\textwidth}@{}}
\toprule
Layer & Meaning in this project\\
\midrule
Code verification & The implementations reproduce their declared discrete
equations, amplification behavior, boundaries, and solve residual.\\
Solution verification & The numerical fields approach the continuous or
semidiscrete reference under a declared refinement.\\
Model validation & Comparison with physical reality or experiments. It is not
performed here.\\
Numerical uncertainty & Refinement, reference evaluation, tridiagonal solve,
and floating-point effects on the reported values.\\
\bottomrule
\end{tabular}
\end{center}

\section{Metrics and observed order}

For interior grid values, the retained study reports
\begin{align}
  E_{\infty}
  &=\max_{1\leq i\leq N-1}
    |\theta_i-\theta_{\mathrm{ref},i}|,\\
  E_{\mathrm{RMS}}
  &=\left[\frac{1}{N-1}\sum_{i=1}^{N-1}
    (\theta_i-\theta_{\mathrm{ref},i})^2\right]^{1/2}.
\end{align}
For a factor-two refinement, the observed order is
\begin{equation}
  p_{\mathrm{obs}}=\frac{\log(E_h/E_{h/2})}{\log 2}.
\end{equation}

\section{Retained study conditions}

All studies end at \(Fo=0.1\). The combined study uses
\(N=10,20,40,80,160\) and \(\mathrm{steps}=N^2/4\), so \(r=0.4\).
The temporal study holds \(N=20\), uses 80, 160, 320, 640, and 1280 steps,
and compares with REF-SEMIDISCRETE-001.

\begin{center}
\begin{tabular}{llrrr}
\toprule
Study & Method & Finest error & Finest order & Value bound\\
\midrule
Combined & FTCS & \num{1.654619e-5} & \num{2.000230} & \num{2.068712e-5}\\
Combined & Crank-Nicolson & \num{1.181738e-5} & \num{1.999746} & \num{1.476826e-5}\\
Temporal & FTCS & \num{1.415673e-4} & \num{1.000467} & \num{1.770737e-4}\\
Temporal & Crank-Nicolson & \num{1.814958e-8} & \num{2.000004} & \num{2.268708e-8}\\
\bottomrule
\end{tabular}
\end{center}

Every error sequence decreases monotonically. The combined results approach
second order. The isolated temporal results approach first order for FTCS and
second order for Crank-Nicolson.

\section{Case-specific acceptance}

For nested fields with target order \(p\), the fine-grid remaining-error
estimate is
\begin{equation}
  U_h=1.25\frac{\|\theta_{h/2}-\theta_h\|_{\infty}}
                    {2^p-1}
      +64\epsilon_{64}\max(1,\|\theta_{h/2}\|_{\infty}).
\end{equation}
The factor 1.25 is the declared safety allowance around the leading-order
Richardson estimate. The second term represents accumulated floating-point
scale. Every finest-grid analytic or semidiscrete error is below this bound.

The order tolerance is twice the absolute change between the two finest
observed orders plus the floating allowance propagated relative to the finest
error. This rule tests entry into the measured asymptotic regime. It replaces
the former unexplained \(\pm0.2\) band.

Direct value comparisons retain the form
\[
  |\mathrm{error}|\leq a_{\mathrm{tol}}
  +r_{\mathrm{tol}}|\mathrm{reference}|,
\]
with terms chosen for the individual comparison. The discrete-mode code tests
use a 64-epsilon relative allowance and an absolute allowance proportional to
the number of time steps.

\section{Numerical uncertainty and limits}

The Richardson term estimates unresolved discretization error for the selected
sequence. Continuous and semidiscrete references are evaluated in IEEE binary64.
The largest retained Crank-Nicolson linear residual is
\(4.44\times10^{-16}\), below its corresponding discretization error by more
than ten orders of magnitude. These observations do not quantify physical
input or model-form uncertainty.

The single smooth Fourier mode does not exercise discontinuities, variable
coefficients, nonlinear sources, nonuniform grids, or multidimensional
operators. The reported orders and tolerances apply only to the declared case.

\end{document}
